\documentclass[10pt]{article}
\usepackage[a4paper,margin=20mm]{geometry}
\usepackage{enumitem}
\usepackage[hidelinks]{hyperref}
\usepackage{microtype}
\setlength{\parindent}{0pt}
\setlength{\parskip}{5pt}
\setlist[enumerate]{leftmargin=*,itemsep=7pt,topsep=4pt}
\begin{document}

\begin{center}
{\Large\bfseries Information Sheet for a Regular Paper}\par
\vspace{4pt}
{\bfseries Same Words, Different Procedure: Counterfactual Transfer Tests for
Reusable LLM-Agent Skills}
\end{center}

\begin{enumerate}
\item \textbf{What is the main claim of the paper, and why is it important to
the autonomous-agents and multi-agent-systems literature?}

Success after supplying an agent with a reusable skill is not, by itself,
evidence that the skill caused success. A base agent may already solve the
target, generic grounding and search controls may explain the gain, or a
lexically similar source may encode an incompatible procedure. The paper's
main claim is that a state-matched, six-condition counterfactual protocol can
separate these explanations and attribute target success specifically to
procedural compatibility under the tested executor.

The Counterfactual Skill Transfer Test (CSTT) crosses compatible versus
incompatible source procedures with near versus far lexical similarity on the
same native target. It adds a no-skill condition and a binding-only condition
that receives the same target grounding and search safeguards but no source
procedure. This distinction matters for agent architecture evaluation: without
it, a skill library can be credited for behavior supplied by the base model or
the executor. CSTT also produces deployment-relevant measures of helpful,
harmful, redundant, and shared-failure transfer. The contribution fits the
journal's interests in evaluating agent decision-making architectures,
testbeds for agent analysis, and software-engineering methods for autonomous
agents.

\item \textbf{What evidence supports the claim?}

After two retained failed pilots and one final mechanism revision, the complete
method, source assignment, three model revisions, inference procedure, decision
gates, and a zero-overlap confirmation roster were frozen. Confirmation uses 60
previously untouched ALFWorld targets, three open local model families, six
conditions per target--model block, 180 paired blocks, and 1,080 native
episodes.

The procedurally matched but lexically far condition succeeds in 176/180 blocks
(97.8\%), compared with 18/180 (10.0\%) for the lexically near procedural lure.
The paired gain is 87.8 percentage points, with a target-stratified 95\%
bootstrap interval of 82.8--93.3 points and a task-clustered randomization
$p=1.00\times10^{-6}$. Relative to the fair binding-only control, matched-far
produces 150 helpful and zero harmful discordances, an 83.3-point net gain. The
primary effect is 88.3, 88.3, and 86.7 points in the three model families and
positive in all six task families. Independent audits accept all 1,080 episodes
and 46,662 native turns with exact within-block initial-state matching, valid
native actions, valid source assignments, and no target-plan leakage.

The claim is deliberately bounded. The confirmation covers instance-level
transfer in one embodied text environment and three small-to-medium open
models. It does not establish universal retrieval quality, cross-domain
transfer, or benefits for proprietary frontier models.

\item \textbf{Which papers make the most closely related contributions, and
how does this paper differ?}

Taylor and Stone, ``Transfer Learning for Reinforcement Learning Domains: A
Survey,'' \emph{Journal of Machine Learning Research} 10 (2009), 1633--1685,
and da Silva and Costa, ``A Survey on Transfer Learning for Multiagent
Reinforcement Learning Systems,'' \emph{Journal of Artificial Intelligence
Research} 64 (2019), 645--703, establish transfer criteria, baselines, and the
importance of source--target relationships. Tuynman and Ortner,
``Quantification of transfer in reinforcement learning via regret bounds for
learning agents,'' \emph{Autonomous Agents and Multi-Agent Systems} 40 (2026),
article 16, quantifies transfer relative to learning without shared experience.
CSTT applies the same comparative principle to reusable LLM-agent procedures,
but uses paired native outcomes and a control that isolates executor assistance.

The closest recent skill studies are Zhou et al., ``Counterfactual Trace
Auditing of LLM Agent Skills'' (arXiv:2605.11946), which labels behavioral
influence patterns in paired with/without-skill traces; Yang et al.,
``SkillMaster'' (arXiv:2605.08693), which uses counterfactual probe utility to
train skill editing; and Belikova et al., ``Managing Procedural Memory in LLM
Agents'' (arXiv:2606.23127), which benchmarks transfer across tasks, roles, and
models. CSTT differs by experimentally crossing procedural compatibility with
lexical similarity on the same initial state and by separating source-procedure
credit from target binding and generic search controls.

\item \textbf{Have parts of this paper been published before?}

No part of this manuscript has been published in an archival journal or
conference, and this manuscript is not under consideration elsewhere. The
authors have released the code, records, and manuscript-supporting artifacts in
a public repository and a versioned Zenodo archive
(\url{https://doi.org/10.5281/zenodo.22852635}) for reproducibility; these
artifact releases are not prior archival publication of the article.

For transparency, a related manuscript by the same authors, ``From Traces to
Witnessed Contracts: Incremental Precondition--Effect Compilation for LLM Agent
Skills,'' is currently under consideration at \emph{Automated Software
Engineering}. That work studies pre-execution detection of defective skill
revisions. The present manuscript studies causal attribution of skill-transfer
benefits. The two studies use disjoint confirmation targets and different
methods, estimands, experiments, figures, tables, endpoints, and results. They
share only the broader agent-skill lifecycle motivation, a public benchmark
platform, and general experimental infrastructure. A copy of the related
manuscript is supplied separately for editorial assessment.
\end{enumerate}

\end{document}
